% Options for packages loaded elsewhere
\PassOptionsToPackage{unicode}{hyperref}
\PassOptionsToPackage{hyphens}{url}
\documentclass[
  10pt,
]{article}
\usepackage{xcolor}
\usepackage[margin=0.8in]{geometry}
\usepackage{amsmath,amssymb}
\setcounter{secnumdepth}{-\maxdimen} % remove section numbering
\usepackage{iftex}
\ifPDFTeX
  \usepackage[T1]{fontenc}
  \usepackage[utf8]{inputenc}
  \usepackage{textcomp} % provide euro and other symbols
\else % if luatex or xetex
  \usepackage{unicode-math} % this also loads fontspec
  \defaultfontfeatures{Scale=MatchLowercase}
  \defaultfontfeatures[\rmfamily]{Ligatures=TeX,Scale=1}
\fi
\usepackage{lmodern}
\ifPDFTeX\else
  % xetex/luatex font selection
\fi
% Use upquote if available, for straight quotes in verbatim environments
\IfFileExists{upquote.sty}{\usepackage{upquote}}{}
\IfFileExists{microtype.sty}{% use microtype if available
  \usepackage[]{microtype}
  \UseMicrotypeSet[protrusion]{basicmath} % disable protrusion for tt fonts
}{}
\makeatletter
\@ifundefined{KOMAClassName}{% if non-KOMA class
  \IfFileExists{parskip.sty}{%
    \usepackage{parskip}
  }{% else
    \setlength{\parindent}{0pt}
    \setlength{\parskip}{6pt plus 2pt minus 1pt}}
}{% if KOMA class
  \KOMAoptions{parskip=half}}
\makeatother
\usepackage{longtable,booktabs,array}
\newcounter{none} % for unnumbered tables
\usepackage{calc} % for calculating minipage widths
% Correct order of tables after \paragraph or \subparagraph
\usepackage{etoolbox}
\makeatletter
\patchcmd\longtable{\par}{\if@noskipsec\mbox{}\fi\par}{}{}
\makeatother
% Allow footnotes in longtable head/foot
\IfFileExists{footnotehyper.sty}{\usepackage{footnotehyper}}{\usepackage{footnote}}
\makesavenoteenv{longtable}
\usepackage{graphicx}
\makeatletter
\newsavebox\pandoc@box
\newcommand*\pandocbounded[1]{% scales image to fit in text height/width
  \sbox\pandoc@box{#1}%
  \Gscale@div\@tempa{\textheight}{\dimexpr\ht\pandoc@box+\dp\pandoc@box\relax}%
  \Gscale@div\@tempb{\linewidth}{\wd\pandoc@box}%
  \ifdim\@tempb\p@<\@tempa\p@\let\@tempa\@tempb\fi% select the smaller of both
  \ifdim\@tempa\p@<\p@\scalebox{\@tempa}{\usebox\pandoc@box}%
  \else\usebox{\pandoc@box}%
  \fi%
}
% Set default figure placement to htbp
\def\fps@figure{htbp}
\makeatother
% definitions for citeproc citations
\NewDocumentCommand\citeproctext{}{}
\NewDocumentCommand\citeproc{mm}{%
  \begingroup\def\citeproctext{#2}\cite{#1}\endgroup}
\makeatletter
 % allow citations to break across lines
 \let\@cite@ofmt\@firstofone
 % avoid brackets around text for \cite:
 \def\@biblabel#1{}
 \def\@cite#1#2{{#1\if@tempswa , #2\fi}}
\makeatother
\newlength{\cslhangindent}
\setlength{\cslhangindent}{1.5em}
\newlength{\csllabelwidth}
\setlength{\csllabelwidth}{3em}
\newenvironment{CSLReferences}[2] % #1 hanging-indent, #2 entry-spacing
 {\begin{list}{}{%
  \setlength{\itemindent}{0pt}
  \setlength{\leftmargin}{0pt}
  \setlength{\parsep}{0pt}
  % turn on hanging indent if param 1 is 1
  \ifodd #1
   \setlength{\leftmargin}{\cslhangindent}
   \setlength{\itemindent}{-1\cslhangindent}
  \fi
  % set entry spacing
  \setlength{\itemsep}{#2\baselineskip}}}
 {\end{list}}
\usepackage{calc}
\newcommand{\CSLBlock}[1]{\hfill\break\parbox[t]{\linewidth}{\strut\ignorespaces#1\strut}}
\newcommand{\CSLLeftMargin}[1]{\parbox[t]{\csllabelwidth}{\strut#1\strut}}
\newcommand{\CSLRightInline}[1]{\parbox[t]{\linewidth - \csllabelwidth}{\strut#1\strut}}
\newcommand{\CSLIndent}[1]{\hspace{\cslhangindent}#1}
\setlength{\emergencystretch}{3em} % prevent overfull lines
\providecommand{\tightlist}{%
  \setlength{\itemsep}{0pt}\setlength{\parskip}{0pt}}
\usepackage{bookmark}
\IfFileExists{xurl.sty}{\usepackage{xurl}}{} % add URL line breaks if available
\urlstyle{same}
\hypersetup{
  hidelinks,
  pdfcreator={LaTeX via pandoc}}

\author{}
\date{}

\begin{document}

\section{Interaction Memory for Reinforcement Learning with Frozen
Vision-Language-Action
Models}\label{interaction-memory-for-reinforcement-learning-with-frozen-vision-language-action-models}

\textbf{Research manuscript, 3 October 2026.} Proposed method and
evaluation protocol; the reported pilot measurements concern the
existing development branch, not a completed evaluation of the proposed
method. Author information is intentionally omitted. This is an extended
working draft, not a submission-formatted or experimentally complete
paper.

\begin{quote}
\textbf{Publication update.} Baseline validation update, 3 October 2026. Before implementing the proposed Stage 1B, validate a reproducible simulation baseline. AlphaBrain on LIBERO is a candidate under source review, not an independently reproduced result. Its released Qwen RLT\_a checkpoint is an evaluation starting point; it differs from the full-token RLT track. Existing ManiSkill pilot measurements below retain their original scope. No new GPU training or control benefit is reported in this publication update.
\end{quote}

\subsection{Abstract}\label{abstract}

Lightweight reinforcement learning can adapt a frozen
vision-language-action model through a compact state representation and
a small actor--critic. However, a current observation need not reveal
how the robot responds to commands under its present execution
conditions. We study whether completed interactions can supply this
missing information. We propose a bounded interaction memory whose
encoder and reader are trained to predict the consequences of subsequent
actions from strictly preceding experience. The resulting context
supplements the original RL token and conditions a small actor and
critic, while the visual-language-action backbone remains frozen during
online learning. Recent ordered records capture immediate execution
history; a bounded archive supports retrieval across verified retries
and, in a separate protocol, transfer from source interactions. The
design separates adaptation through memory updates from learning through
parameter updates. Existing simulation diagnostics show that
command--response history contains predictive information, but also
expose a weakness: a fixed response estimator substantially outperforms
the current attention reader. These observations motivate
consequence-supervised memory training without establishing a control
benefit. We specify controlled experiments for online sample efficiency,
frozen-parameter adaptation, stale-memory effects, visual relevance, and
transfer. Closed-loop improvements and reduced human effort remain
hypotheses to be tested.

\subsection{1. Introduction}\label{introduction}

This work investigates how a robot can use completed physical
interactions to improve lightweight reinforcement learning around a
frozen vision-language-action (VLA) model. A pretrained VLA supplies
useful task and visual information, but two observations with similar
appearance and joint configuration can precede different responses to
the same command. A robot may encounter a different controller gain,
delayed execution, an object that slips, or contact that blocks motion.
Some of these differences become observable only after acting. The
relevant question is whether the agent can retain that evidence and use
it when choosing its next action.

RLT provides a suitable experimental starting point: a compact token
connects a VLA to lightweight online actor--critic learning (Xu et al.
2026). Its computational advantage makes controlled experiments feasible
on modest hardware. Our local RLT implementation also exposes a
reference action chunk and proprioception to the small controller. We
preserve these inputs. We investigate an additional source of
information: what actually happened when commands were executed. The
objective is not to replace the visual token or fine-tune a large
backbone online.

There are two distinct ways that experience can affect behavior.
Replay-based reinforcement learning changes network parameters using
past transitions. A context memory changes the information available to
a decision, even when parameters remain fixed. These mechanisms need
separate evaluation. Ordinary RL already uses unsuccessful transitions
through reward and value learning; storing a failure is not a new
algorithmic contribution. A memory mechanism is useful only if its
contents change predictions or decisions in a beneficial and
reproducible way.

Our proposal is inspired by Zeva's use of completed interactions and
deployment-time memory (Chen et al. 2026). We ask a narrower, testable
question: \textbf{Can a representation learned to extract future-useful
execution evidence improve an RLT controller under a fixed interaction
and compute budget?} ``Future-useful'' has an operational definition.
The memory encoder must help predict the effect of a later action from
earlier records; reconstructing the effect already stored in the same
record is insufficient evidence. This training criterion gives the
memory a role beyond generic attention over a trajectory.

The proposed contribution is a method and an evaluation of
consequence-supervised context for frozen-VLA RL. A prediction loss plus
concatenated memory is not sufficient novelty by itself; this draft does
not establish that the proposed combination clears that bar. Its
intended empirical contributions are (i) improved online learning at
matched data cost, (ii) useful adaptation from retained evidence with
parameters frozen, and (iii) selective reuse of source experience
without systematic negative transfer. None is asserted as an established
result in this draft. Context inference, response prediction, memory
conditioning, and residual-based reliability have substantial
precedents. The scientific contribution must therefore be demonstrated
by the controlled combination, its failure analysis, and its benefits
over strong history and representation baselines. If the result reduces
to a standard Zeva-style encoder attached to RLT without an additional
empirical finding, substantial novelty remains unresolved.

\subsection{2. Related Work}\label{related-work}

\subsubsection{2.1. Lightweight reinforcement learning around pretrained
action
models}\label{lightweight-reinforcement-learning-around-pretrained-action-models}

VLA models such as pi0 and pi0.5 combine pretrained visual-language
representations with action generation (Black et al. 2024; Intelligence
et al. 2025). RLT studies compact representations for lightweight online
RL (Xu et al. 2026). DSRL instead learns to steer a frozen diffusion
policy through its sampling variables (Wagenmaker et al. 2025). ConRFT
studies offline and online reinforced VLA fine-tuning using a
consistency policy (Chen et al. 2025). These approaches differ in which
parameters and action interfaces are optimized; published success
numbers cannot be compared directly across their training data and
control budgets.

The recently released eRLT is a particularly close representation
baseline. It selects action-relevant information across tokens and
layers of a frozen VLA, initializes routing with expert-action
prediction, and adapts routing through the critic (Huang et al. 2026).
Our proposed information source is completed execution history.
Current-observation routing and history-based condition inference are
complementary hypotheses, but a stronger current-state representation
could eliminate an apparent memory advantage. Accordingly, an eRLT-style
representation comparison is required before claiming that memory
addresses a limitation specific to frozen VLA features.

\subsubsection{2.2. Memory and adaptation from
interaction}\label{memory-and-adaptation-from-interaction}

Zeva learns an interaction representation with action-conditioned
pre-interaction and post-interaction effect prediction, then uses
short-term and cross-attempt memory to condition action generation; its
deployment procedure updates memory with neural parameters frozen (Chen
et al. 2026). We borrow that interaction--storage--conditioning pattern.
Our implementation conditions RLT's small actor and critic and
explicitly separates online weight learning from frozen-parameter memory
adaptation. Effect prediction and dual-timescale memory are inherited
ideas, not new contributions here. A simplified reader inside RLT must
be labeled a Zeva-inspired adaptation rather than a reproduction of the
complete system.

MemoryVLA retrieves perceptual and cognitive information to condition
manipulation (Shi et al. 2025). RoboMME provides a benchmark for memory
in robotic generalist policies (Dai et al. 2026). These works make a
generic ``adding memory improves manipulation'' claim inadequate. Our
evaluation isolates executed commands and responses, distinguishes
visual context from joint history, and tests transfer and interference
under changes in execution conditions.

The broader adaptation literature already formalizes history-conditioned
control. UP-OSI predicts dynamics parameters from recent states and
actions (Yu et al. 2017); RMA learns rapid adaptation from recent
experience (Kumar et al. 2021). PEARL infers context for off-policy
meta-RL, VariBAD conditions behavior on inferred task uncertainty, and
RL2 learns an adaptation procedure through recurrent state (Rakelly et
al. 2019; Zintgraf et al. 2019; Duan et al. 2016). We do not claim the
first physical context representation. Our setting differs through
frozen VLA features, bounded explicit evidence, action-chunk execution,
and the need to measure benefits under a small online budget.

\subsubsection{2.3. Predictive representations and failure
reuse}\label{predictive-representations-and-failure-reuse}

SPR learns representations using future latent prediction (Schwarzer et
al. 2020). FLARE integrates future latent modeling with action
generation (Zheng et al. 2025). Imagine-RL augments VLA reinforcement
learning with a frozen world model and residual-confidence-guided critic
fusion (Hu et al. 2026). TD-MPC2 uses learned latent dynamics for
trajectory optimization (Hansen et al. 2023). We use prediction to train
a representation of preceding interactions; the proposed actor does not
query a predictor for each candidate action or perform model-based
planning. A prediction-only sidecar is nevertheless a necessary control
for separating memory from an auxiliary-loss benefit.

Off-policy actor--critic methods already learn from unfavorable
transitions (Haarnoja et al. 2018; Fujimoto et al. 2018). Hindsight
Experience Replay additionally changes the goal interpretation of failed
attempts (Andrychowicz et al. 2017). HIL-SERL demonstrates a
human-in-the-loop learning system where demonstrations and corrections
matter (Luo et al. 2024). Our failed interactions contribute to dynamics
supervision and value learning, but are not automatically
behavior-cloned. Simulation correction counts are proxies; human minutes
saved must eventually be measured on physical hardware.

\textbf{Table 1. Closest conceptual comparisons.} Entries describe the
mechanisms relevant to this proposal, not exhaustive capabilities of
each work.

{\def\LTcaptype{none} % do not increment counter
\begin{longtable}[]{@{}
  >{\raggedright\arraybackslash}p{(\linewidth - 6\tabcolsep) * \real{0.2500}}
  >{\raggedright\arraybackslash}p{(\linewidth - 6\tabcolsep) * \real{0.2500}}
  >{\raggedright\arraybackslash}p{(\linewidth - 6\tabcolsep) * \real{0.2500}}
  >{\raggedright\arraybackslash}p{(\linewidth - 6\tabcolsep) * \real{0.2500}}@{}}
\toprule\noalign{}
\begin{minipage}[b]{\linewidth}\raggedright
Work
\end{minipage} & \begin{minipage}[b]{\linewidth}\raggedright
Information central to the comparison
\end{minipage} & \begin{minipage}[b]{\linewidth}\raggedright
Learned or updated component relevant here
\end{minipage} & \begin{minipage}[b]{\linewidth}\raggedright
Required distinction
\end{minipage} \\
\midrule\noalign{}
\endhead
\bottomrule\noalign{}
\endlastfoot
RLT & Current VLA representation & Compact representation, then small RL
controller & Preserve baseline action and routing contract \\
eRLT & Tokens and layers of current VLA input & Action-initialized,
critic-refined routing & Test whether better current features remove the
gain \\
Zeva & Completed visual-action interactions & Interaction encoder;
deployment memory & Separate context adaptation from RL weight
updates \\
MemoryVLA & Retrieved perceptual/cognitive history & Memory-conditioned
action model & Compare with generic visual history \\
RMA / UP-OSI & Recent state-action evidence & Adaptation or
identification module & Include a strong recurrent history baseline \\
FLARE / SPR & Future representations & Predictive representation
learning & Include prediction without an archive \\
Imagine-RL & Predicted futures and model residuals &
World-model-augmented value estimation & Do not claim novelty for
predictive critic confidence \\
Proposed method & Earlier execution records useful for later
consequences & Stage 1B context encoder; Stage 2 actor--critic &
Establish benefit through memory interventions \\
\end{longtable}
}

\subsection{3. Problem Formulation}\label{problem-formulation}

We consider partially observed manipulation with visual observations,
instructions, and proprioception. Let the hidden physical state be
\texttt{x\_t}, and let \texttt{theta\_t} summarize execution conditions
such as controller settings, friction, or delay. The transition law
depends on both. \texttt{theta\_t} is a modeling device; its true value
is not provided to the proposed controller and need not be identifiable.
We initially study stationary conditions within an attempt and
subsequently introduce controlled changes.

The frozen Stage 1 model produces an RL token \texttt{z\_t}, transformed
proprioception \texttt{p\_t}, and a reference action chunk
\texttt{A\_ref,t}. Raw joint positions \texttt{q\_t} are separately
available to the interaction collector. Their units differ from the
transformed VLA state. We denote the deployable current input by
\texttt{s\_t\ =\ {[}z\_t,\ p\_t,\ q\_t{]}}; the no-memory reference
receives the same current observations wherever a comparison requires
input matching. Any baseline that gains a new raw-state input is
separately named.

At decision time, the controller proposes a chunk of at most
\texttt{H\ =\ 10} normalized commands. Existing routing selects the
executed commands, which may originate from the VLA, actor, or an
explicitly enabled corrective controller. We store the actual executed
prefix of length \texttt{ell\_t}, not an unselected proposal. The
available history contains only completed records:

\[\mathcal{H}_t = \{(s_i,A_i^{\mathrm{exec}},s_{i+1},m_i):i<t\}. \tag{1}\]

The mask \texttt{m\_i} identifies valid commands, terminal observations,
and stream boundaries. The memory context \texttt{c\_t} is a
deterministic function of a bounded subset of this history and the
current observation. It is an inferred condition, not a calibrated
posterior over physical parameters. We optimize expected discounted task
return while tracking interaction count, online compute, total
pretraining cost, and memory overhead.

A \textbf{physical interaction representation} is defined operationally:
it summarizes how commands and observed outcomes relate in a way that
improves held-out consequence prediction or control. This does not imply
discovery of a unique friction coefficient, contact force, or universal
physical law. Small joint motion can result from obstruction, controller
gain, saturation, or latency. Claims of contact recognition require
independent contact measurements; claims of causal identification
require stronger assumptions or interventions than sequence prediction
alone.

\subsection{4. Consequence-Supervised Interaction
Memory}\label{consequence-supervised-interaction-memory}

\begin{figure}
\centering
\pandocbounded{\includegraphics[keepaspectratio,alt={Figure 1. Proposed information flow. The observation and action backbone is frozen during Stage 1B and Stage 2. Blue blocks are newly proposed context modules. Memory becomes available only after execution. Stage 1B trains context to predict a subsequent outcome; Stage 2 initially freezes that context interface and trains the small actor and critic.}]{figures/framework.png}}
\caption{Figure 1. Proposed information flow. The observation and action
backbone is frozen during Stage 1B and Stage 2. Blue blocks are newly
proposed context modules. Memory becomes available only after execution.
Stage 1B trains context to predict a subsequent outcome; Stage 2
initially freezes that context interface and trains the small actor and
critic.}
\end{figure}

\subsubsection{4.1. Encode completed
events}\label{encode-completed-events}

For each completed chunk, the collector constructs a record containing
the raw start state, executed commands, valid duration, and observed
effects. The current branch already stores raw joint positions and their
changes. We propose adding current and successor visual tokens to
distinguish interactions that have similar joint configurations but
different object arrangements:

\[e_i=E_{\psi}(z_i,q_i,A_i^{\mathrm{exec}},\Delta q_i,z_{i+1},m_i). \tag{2}\]

\texttt{E\_psi} is a small event encoder with a 64-dimensional output. A
two-layer temporal attention block encodes the ordered commands and
their masks; an MLP projects the remaining inputs. This preserves action
order, which a sum-of-commands response formula discards. Initial widths
are 128 hidden units and four attention heads; these are starting
configurations rather than established optimal hyperparameters.
Train-only normalization is stored with the encoder version. No
additional VLA forward pass is required if tokens are already computed
for current and true successor observations.

A record may contain its own observed effect because it is written after
that effect occurs. It may not contain the effect of the action
currently being selected. Termination flags are needed for transport and
validity, but the default learned physics input excludes success and
task reward. Reward remains in replay for value learning. An ablation
includes end flags to test shortcut learning. Invalid padded values must
be removed before an MLP, normalization, attention, or nonlinear loss;
multiplying a NaN-containing scalar by zero is insufficient.

\subsubsection{4.2. Read short-term and longer-term
evidence}\label{read-short-term-and-longer-term-evidence}

The short store contains the latest four valid records in chronological
order. The archive contains up to 32 completed records, including
records that may also be recent; the reader selects up to four older
candidates after excluding the short store. We retain individual records
instead of averaging contradictory outcomes. The first implementation
uses a bounded FIFO archive to isolate representation learning from a
sophisticated storage policy. ``Long-term'' here means older than the
immediate context window; persistence across attempts is an explicit
protocol, not the default meaning of the word.

A query from current visual and proprioceptive inputs attends to the
recent records and selected older records:

\[c_t=R_{\rho}(z_t,q_t,\{e_i:i\in\mathcal{S}_t\cup\mathcal{K}_t\}). \tag{3}\]

For an implementable first learned retrieval scheme, use normalized
current/start visual projections and scaled joint distances to select
archive candidates. Train the visual projection jointly with the context
loss using soft attention over the full 32-record archive in Stage 1B;
apply top-four selection and fine-tune with the deployment selection
rule before freezing. Report the soft-to-hard gap. A simpler fixed
visual-plus-joint distance reader is the first fallback and a baseline.
The reader receives temporal age and source-domain tags that are
observable at deployment, never hidden stiffness or target-test task
labels. Action and effect content are encoded in values; retrieval is
based on available current context.

Empty memory maps to exactly zero. A fixed, validation-selected
context-dropout rate during training exposes the controller to empty
memory, reducing a train/test mismatch in clearing interventions. Report
both natural-retention performance and intervention sensitivity. A
short-only version retains the same output dimension. Cross-attempt
retention requires a verified reset instance and unchanged physical
configuration in the retry experiment. Cross-task transfer uses a
separate source bank with declared provenance, compatible
robot/controller units, and an encoder version. The source bank is
read-only during held-out target evaluation; target evidence lives in a
separate temporary bank and is discarded between instances.

\subsubsection{4.3. Train memory for a later
outcome}\label{train-memory-for-a-later-outcome}

Stage 1B trains the event encoder and reader on a fixed transition
dataset after the original Stage 1 checkpoint has been selected. It
reconstructs the history available immediately before each target
action. The decoder sees the current observation, this past context, and
the target action sequence. It predicts the target action's subsequent
effect:

\[ (\widehat{\Delta q}_t,\hat z_{t+1})=D_{\omega}(z_t,q_t,c_t,A_t^{\mathrm{exec}},m_t). \tag{4}\]

The initial predictive dataset uses fixed-horizon windows with valid
successors and reports any terminal-window exclusions. Outcome-dependent
termination length and success flags are not prediction inputs. Online
residual checks commit to a horizon before execution and bind feedback
to that prediction. All context records in Equation (4) precede the
target chunk. The target outcome is excluded from context construction,
key selection, normalization fitting, and memory summaries. Windows from
one underlying trajectory or paired scene/command seed stay in the same
dataset split. The action input is available to this training objective
but the prediction head is not required for action selection.

We supervise raw joint changes in normalized coordinates and a fixed
visual target. Direct prediction is the initial visual head: a separate
development branch found that residual latent prediction was not
uniformly better at long horizons. We retain both as an ablation rather
than imposing a residual parameterization:

\[\mathcal{L}_{\mathrm{ctx}}=\lambda_q\operatorname{Huber}(\widehat{\Delta q}_t,\Delta q_t)+\lambda_z[1-\cos(\hat z_{t+1},\operatorname{sg}(z_{t+1}))]. \tag{5}\]

Start with \texttt{lambda\_q\ =\ 1} after train-split standardization
and choose \texttt{lambda\_z} from \texttt{\{0,\ 0.1,\ 1\}} using
validation episodes only. The frozen token target prevents target drift.
The loss does not ensure memory use: the decoder might solve the problem
from the current input alone. We therefore require no-history,
shuffled-history, command-shuffled, and mismatched-context comparisons
on the same held-out groups. Those interventions diagnose dependence;
they do not establish physical identifiability.

Demonstrations often cover successful execution but few blocked or
delayed responses. Stage 1B therefore uses a shared, explicitly costed
pool of demonstrations and bounded exploratory interactions, including
failures. Every comparison receives the same source trajectories;
prediction-free baselines may use them for their permitted RL warm
start. We report both online-only sample efficiency and total-data
efficiency. A benefit purchased with additional private pretraining data
is not an equal-data gain.

An optional contact head can predict a simulator contact label from
deployable inputs during training. This is privileged supervision and
must be labeled accordingly, counted in the comparison, and ablated. The
core proposal uses visual and joint effects alone. Desired task contact
is not universally penalized; a contact-count reduction is meaningful
only alongside success and a task-specific definition of undesirable
impact.

\subsubsection{4.4. Condition RLT without destabilizing the
representation}\label{condition-rlt-without-destabilizing-the-representation}

The first proposed Stage 2 freezes the VLA, original token encoder, and
pretrained event encoder/reader. It trains only the small actor and twin
critics. This departs deliberately from the current branch, whose reader
learns through TD loss. Freezing makes the context coordinates and
transferred bank stable while testing whether the representation
contains useful information. A later ablation jointly updates the
context modules with TD and the auxiliary loss; it must re-encode raw
replay snapshots and version any cached memory keys.

The actor input is \texttt{{[}A\_ref,t,\ z\_t,\ p\_t,\ c\_t{]}}; the
critic input is \texttt{{[}z\_t,\ p\_t,\ c\_t,\ A\_t{]}}. If raw
\texttt{q\_t} is also exposed directly, the capacity-matched baseline
receives it too. The original routing, action bounds, reward,
demonstration sampling, exploration, and update ratio remain fixed
within each comparison. The memory reader never owns a second optimizer
through the actor path.

With \texttt{R\_t} the discounted reward inside a chunk and
\texttt{b\_t} the baseline bootstrap-validity mask, the intended
semi-Markov target is

\[y_t=R_t+\gamma^{\ell_t} b_t\min_{j=1,2}\bar Q_j(s_{t+1},c_{t+1},\bar\pi(s_{t+1},c_{t+1})). \tag{6}\]

The baseline implementation's fixed reward horizon and
termination/truncation handling must be preserved or corrected uniformly
across all methods before comparison; using one discount factor per
ten-step chunk changes the algorithm. For transitions that terminate
early, target construction must respect the actual valid execution and
the environment's terminal-successor contract.

\[\mathcal{L}_Q=\sum_{j=1}^{2}\mathbb{E}_{\mathcal{B}}[(Q_j(s_t,c_t,A_t^{\mathrm{exec}})-\operatorname{sg}(y_t))^2]. \tag{7}\]

The actor keeps the baseline value objective plus behavior-cloning terms
on reference or correction targets. Failed executed actions are useful
for Equation (7) and the consequence loss, but do not become positive
imitation targets merely because they are stored. Entropy and
action-noise settings remain those of the matched baseline; this draft
does not introduce a new RL optimizer.

\subsubsection{4.5. Maintain temporal integrity and handle stale
experience}\label{maintain-temporal-integrity-and-handle-stale-experience}

Replay stores owned copies of the context evidence visible at both
current and real successor decisions. Recomputing an old transition
against today's archive would leak future information. The event from
the current action enters \texttt{c\_next} only after its true successor
is observed. It never enters \texttt{c\_t}. Partial resets affect only
their simulator lanes. A process restart does not silently imply that
the environment archive has been restored.

The core method uses a fixed memory policy. A separate robustness
variant uses observed prediction errors to downweight or reset older
evidence after a validated change signal. Its threshold is selected on
validation streams, with a stationary-condition false-alarm constraint.
Excitation support is not predictive confidence: abundant old evidence
can be wrong after a condition change. Delayed feedback is matched to
the prediction and memory generation that produced it; feedback from
before a reset cannot repeatedly reset the new generation. This variant
is evaluated in actual changed-condition simulation, not only synthetic
scalar response streams. It is not presented as a new principle of
confidence estimation.

\textbf{Algorithm 1. One online interaction cycle.}

\begin{enumerate}
\def\labelenumi{\arabic{enumi}.}
\tightlist
\item
  Read current visual/proprioceptive inputs and retrieve strictly
  preceding, permitted memory records. Compute \texttt{c\_t}.
\item
  Sample a small-actor action and apply the unchanged RLT routing rule.
\item
  Execute the valid command prefix and collect rewards and the true
  pre-reset successor.
\item
  Encode the completed event, append it to the permitted stores, and
  construct successor evidence.
\item
  Store executed actions, valid duration, terminal masks, and immutable
  current/next evidence in replay.
\item
  Update actor and critics under the matched RL schedule. Keep
  representation modules frozen in the primary configuration.
\item
  On reset, apply the declared clear/retain protocol. At
  frozen-parameter evaluation, skip all gradient updates.
\end{enumerate}

\subsubsection{4.6. Why action-expert features are an
extension}\label{why-action-expert-features-are-an-extension}

The action expert may expose execution-relevant internal information,
but it introduces a second representation intervention. During local
Stage 1 training, its hidden states depend on noisy ground-truth action
inputs and flow time. Those states are not automatically available under
the same distribution at deployment. A valid experiment must specify a
reproducible inference-time extraction rule, such as a fixed denoising
time and fixed noise convention applied to a reference action proposal,
and measure its extra latency. We defer this extension until history
conditioning demonstrates a benefit. Joint Stage 1 training can
subsequently add the consequence objective while preserving token
reconstruction and action losses; it must use actual successor-aligned
trajectories, not independent observation-action examples with invented
next states.

\subsection{5. Evaluation Design}\label{evaluation-design}

\subsubsection{5.1. Questions and preregistered
outcomes}\label{questions-and-preregistered-outcomes}

The main hypothesis is that completed interaction evidence improves
normalized success-rate AUC over online environment transitions under
matched source data and updates. Final success, successful episodes per
wall-clock hour, and total collection cost are co-reported. A lower
prediction error is a diagnostic outcome, not the primary control
result. Four additional questions test mechanism: does retaining memory
help with weights frozen; does visual context improve over joint-only
history; can outdated records harm adaptation; and can a source bank
help a new condition or task beyond transferred network weights?

Before confirmation experiments, freeze the task list, seeds, primary
metric, evaluation states, update budget, checkpoint rule, and source
data. Use the final-budget checkpoint for the primary result; any
validation-selected checkpoint is a separate report. No checkpoint is
selected from test success. An initially proposed non-inferiority margin
is five percentage points of final success; it is only claimed when the
lower confidence bound of the paired difference exceeds minus that
margin.

\subsubsection{5.2. Benchmarks and condition
splits}\label{benchmarks-and-condition-splits}

We start with the locally integrated ManiSkill peg-insertion task (Tao
et al. 2024). The existing wide geometry is an engineering starting
point; narrower geometry is a held-out geometry condition, not an
independent task family. Add a pushing/alignment task and an articulated
manipulation task once their demonstrations, action controller, and
visual Stage 1 checkpoint are validated. PushCube and TurnFaucet are
candidates, not claims of already supported local configurations. A
second benchmark such as LIBERO can test broader transfer later (Liu et
al. 2023), but switching simulator and controller prematurely would
confound the first result.

Training varies supported physical parameters within a declared range.
Separate tests hold out (a) scene/command seeds, (b) physical parameter
values, (c) combinations of known parameter values, (d) geometry or
object instances, and (e) task identity. These are different
generalization claims. Controller gain, execution delay, friction, and
object mass are changed one at a time initially. A wrapper that clips
commands is an execution perturbation, not a simulated friction change.
Read back the effective controller/physics parameters after reset and
verify that the intended intervention actually took effect.

At least one task should require visual disambiguation: comparable robot
configurations and command histories with different object alignment or
contact geometry. Joint-only baselines may perform well on
controller-gain identification; that does not establish a visual
reasoning contribution. If the visual event features are unnecessary
throughout, the eventual paper should narrow its claims and reconsider
its fit to a computer-vision venue.

\subsubsection{5.3. Baselines and controlled
budgets}\label{baselines-and-controlled-budgets}

\textbf{Table 2. Main comparison plan.} Original RLT remains a reference
even if a wider no-memory controller is used for capacity control. All
memory methods use the same source trajectories and online interaction
budget.

{\def\LTcaptype{none} % do not increment counter
\begin{longtable}[]{@{}
  >{\raggedright\arraybackslash}p{(\linewidth - 6\tabcolsep) * \real{0.2500}}
  >{\raggedright\arraybackslash}p{(\linewidth - 6\tabcolsep) * \real{0.2500}}
  >{\raggedright\arraybackslash}p{(\linewidth - 6\tabcolsep) * \real{0.2500}}
  >{\raggedright\arraybackslash}p{(\linewidth - 6\tabcolsep) * \real{0.2500}}@{}}
\toprule\noalign{}
\begin{minipage}[b]{\linewidth}\raggedright
ID
\end{minipage} & \begin{minipage}[b]{\linewidth}\raggedright
Method
\end{minipage} & \begin{minipage}[b]{\linewidth}\raggedright
Context training
\end{minipage} & \begin{minipage}[b]{\linewidth}\raggedright
Primary purpose
\end{minipage} \\
\midrule\noalign{}
\endhead
\bottomrule\noalign{}
\endlastfoot
B0 & Original local RLT & None & Preserve the starting baseline \\
B1 & Capacity-matched RLT, zero context & None & Separate input width
and parameter capacity \\
B2 & Recent ordered visual/proprioceptive history & TD-only recurrent
reader & Strong history and partial-observability control \\
B3 & Current short-plus-archive reader & TD-only & Measure the existing
Zeva-inspired branch \\
B4 & Fixed response features, padded to 64 dimensions & Parameter-free
reader & Test the strong diagnostic alternative in control \\
M & Proposed consequence-supervised memory & Stage 1B, then frozen in
primary Stage 2 & Test future-useful history conditioning \\
N1 & eRLT-style current representation & Action initialization and TD
adaptation & Nearest representation comparison \\
N2 & Prediction-only sidecar, no archive & Matched consequence
supervision & Isolate the auxiliary-prediction effect \\
N3 & Zeva-style CTE feeding the same RLT controller & Pre/post effect
objectives and matched history & Test against the closest
interaction-encoding design \\
\end{longtable}
}

N1 must either reproduce a released compatible implementation or be
labeled an adaptation with architecture and cost differences. DSRL and
full Zeva are useful external comparisons only if task, source data,
backbone, and intervention budgets can be aligned. Otherwise, discuss
them without placing incomparable paper numbers in a ranking table. At
minimum, N1 or a clearly justified action-supervised current-token
control is required before a broad representation claim. N3 is a
targeted mechanism control: compare the same history/data budget and
downstream controller, disclose omitted Zeva losses or retrieval
components, and do not present an incomplete adaptation as a full
reproduction.

Stage 1B cost is charged explicitly. Report both fixed online
interaction and fixed end-to-end compute comparisons. Match
update-to-data ratios, BC pretraining, replay capacities, number of
environments, action frequency, chunk size, and evaluation overhead.
Actor initialization shares common weights where tensor shapes permit;
added inputs are identically initialized across capacity controls. If
the existing RLT phase gate uses simulator state, disclose it and keep
it identical. ``No corrective expert'' does not mean ``no privileged
routing information.'' Always-on actor experiments form a separate
controlled block.

\subsubsection{5.4. Ablations that can reject the
explanation}\label{ablations-that-can-reject-the-explanation}

\textbf{Table 3. Mechanism tests, ordered by scientific priority.}

{\def\LTcaptype{none} % do not increment counter
\begin{longtable}[]{@{}
  >{\raggedright\arraybackslash}p{(\linewidth - 4\tabcolsep) * \real{0.3333}}
  >{\raggedright\arraybackslash}p{(\linewidth - 4\tabcolsep) * \real{0.3333}}
  >{\raggedright\arraybackslash}p{(\linewidth - 4\tabcolsep) * \real{0.3333}}@{}}
\toprule\noalign{}
\begin{minipage}[b]{\linewidth}\raggedright
Intervention
\end{minipage} & \begin{minipage}[b]{\linewidth}\raggedright
What is kept fixed
\end{minipage} & \begin{minipage}[b]{\linewidth}\raggedright
What a negative result means
\end{minipage} \\
\midrule\noalign{}
\endhead
\bottomrule\noalign{}
\endlastfoot
Remove history at train and test & Data, context width, decoder capacity
& Current inputs may suffice \\
Remove commands from event; remove target action separately and jointly
& All observation windows & Joint removal tests action dependence;
separate removals locate its source \\
Shuffle history within matched strata & Marginal records, count,
normalization & Reader may ignore event correspondence \\
Recent only versus recent plus archive & Context width and training &
Archive may be unnecessary \\
No Stage 1B versus consequence supervision & Source data access and
Stage 2 budget & Auxiliary training may not help control \\
Visual-plus-joint versus joint-only events & Current visual input,
history budget & Visual memory contribution may be absent \\
Predict only joints versus joints plus visual effects & Module capacity
& Visual target could be nuisance supervision \\
Direct versus residual future latent head & Data, loss scale, updates &
Residual structure may be harmful \\
Actor-only versus critic-only versus both conditioning & Shared learned
context & Locate where information helps \\
Frozen versus online-updated context modules & Initialization, source
data & Quantify stability/adaptability tradeoff \\
Uniform/FIFO versus learned retrieval & Record budget and reader &
Learned retrieval may not justify complexity \\
Retain/clear/shuffle at evaluation & Identical trained weights and
initial states & Isolate benefit of the evidence itself \\
\end{longtable}
}

Do not train a separate model on every optional combination before the
core causal comparisons succeed. The archive, reliability variant, and
action-expert extension enter sequentially. A negative ablation is not
hidden by averaging it with easier tasks.

\subsubsection{5.5. Frozen-parameter retries and
transfer}\label{frozen-parameter-retries-and-transfer}

For retry adaptation, fix every network parameter and optimizer state.
Run five attempts for each held-out instance with identical reset state
and physical conditions. Compare retaining the archive, clearing it, and
supplying a count-matched shuffled archive. Continue the attempt
schedule after success, resetting the same instance, to avoid selecting
only hard survivors. Report unconditional success at each attempt,
success conditional on previous failure as a distinct quantity, and
cumulative success. Even an unchanged independent controller achieves
cumulative success \texttt{1\ -\ (1\ -\ p)\^{}K}; increasing cumulative
success alone does not establish learning.

For transfer, hold the target initialization and weights fixed while
changing only the source bank. Compare empty bank, related source
interactions, unrelated but count-matched interactions, and a bank from
a different physical condition. Separately compare pretrained versus
untrained context weights without a source bank. This separates
knowledge stored in parameters from knowledge stored in records.
Target-test trajectories never enter the source bank or hyperparameter
selection. Cross-task transfer first uses the same robot, action
convention, and sensors; a new embodiment requires a separate alignment
method and is outside the first claim.

Changed-condition tests retain an old bank deliberately, switch a
verified physical parameter, and compare immediate clearing, no
clearing, fixed forgetting, and residual-triggered downweighting. Report
initial damage, recovery interactions, stationary false alarms, and
recovery success. An oracle change indicator is only an upper bound.
Successful adaptation should not depend on access to the hidden switch
time.

\subsubsection{5.6. Statistical and efficiency
reporting}\label{statistical-and-efficiency-reporting}

The pilot uses three training seeds and small validation panels to find
implementation errors. Confirmation should use five independent training
seeds for the core baseline--method contrast when feasible, and at least
three for secondary ablations. Report seed-level results, not only a
pooled mean. Final evaluation uses at least 100 held-out episodes per
seed and condition, increased if a non-inferiority or small-effect claim
requires tighter precision. For intuition, a success probability near
0.5 has a roughly ten-point binomial 95\% half-width with only 100
independent episodes; training variability adds uncertainty.

Use paired reset seeds and controlled action-noise streams where
applicable. Bootstrap hierarchically over training seeds and evaluation
instances; never treat adjacent trajectory windows or retries as
independent training runs. Report uncertainty as recommended by the RL
evaluation literature (Agarwal et al. 2021), while acknowledging that
three seed clusters provide weak interval estimates. Predefine the main
contrast and apply a stated multiplicity correction to families of
confirmatory secondary tests.

Normalize learning-curve AUC by the fixed interaction budget so its
scale remains success probability. Count valid environment control ticks
as well as action chunks and optimizer steps. Threshold-hitting time is
censored when a run never reaches the target; do not average only
successful runs. Report simulator FPS, inference latency percentiles,
end-to-end environment steps per second, successful episodes per hour,
GPU-hours, peak memory, and archive/replay bytes. A faster success curve
per transition can coexist with worse wall-clock performance.

Corrections, resets, and collision proxies are secondary simulation
measurements with a fixed trigger and fixed corrective-controller budget
across methods. Human correction time, setup time, and recovery time
require real operator measurements and remain future work.

\subsection{6. Existing Development
Evidence}\label{existing-development-evidence}

The following measurements were collected before the proposed method was
implemented. They motivate its design and delimit current claims. They
are not held-out confirmation results for the new representation.

\subsubsection{6.1. Response prediction from physical
simulation}\label{response-prediction-from-physical-simulation}

The existing hidden-PD diagnostic contains 56 paired scene/command
groups, each executed under stiffness 250 and 1000, for 112 trajectories
of 120 control ticks. The 32/8/16 group
train/validation/development-test split keeps each pair together. Both
stiffness values occur in every split. Ten-step windows yield
768/192/384 examples, but these windows are correlated and are not
independent environments. The learned probes use three initialization
seeds and 1,000 updates, with checkpoint selection on validation groups.

\textbf{Table 4. Observed development-test joint-change MSE. Lower is
better.} Means for learned methods; no uncertainty interval was
reconstructed here. The fixed formula has structural knowledge and is
not parameter matched.

{\def\LTcaptype{none} % do not increment counter
\begin{longtable}[]{@{}lr@{}}
\toprule\noalign{}
Existing reader/probe & MSE \\
\midrule\noalign{}
\endhead
\bottomrule\noalign{}
\endlastfoot
No history & 1.8211e-4 \\
Recent history & 1.8176e-4 \\
Archive only & 1.8267e-4 \\
Full attention history & 1.8143e-4 \\
Response features plus learned prediction head & 1.7533e-4 \\
Fixed command--response formula & 7.3889e-5 \\
\end{longtable}
}

The attention result is close to the no-history result, whereas a
structured response estimator is substantially better. This supports the
narrower inference that response-relevant information exists in these
records and is poorly extracted by the current neural reader. It does
not show improved task success or extrapolation to unseen stiffness. A
fresh confirmation split is required because the development set has
already influenced design decisions.

These legacy errors aggregate raw joint coordinates, including revolute
arm joints and prismatic fingers; they should not be interpreted as a
single physical unit such as squared radians. New experiments report arm
and finger errors separately and provide train-normalized aggregate
errors. The fixed estimator computes a per-joint command-to-change slope
and an excitation support statistic. It sums normalized arm commands
after applying the known controller scale; it does not estimate force or
identify a unique physical cause. It also loses within-chunk action
order and can be unreliable when opposing commands cancel.

\begin{figure}
\centering
\pandocbounded{\includegraphics[keepaspectratio,alt={Figure 2. Observed development errors. Left: existing reader comparison. Right: a separate post-hoc follow-up. Means are shown without confidence intervals; the fixed formula is not parameter matched. These are prediction diagnostics, not closed-loop results.}]{figures/pilot_diagnostics.png}}
\caption{Figure 2. Observed development errors. Left: existing reader
comparison. Right: a separate post-hoc follow-up. Means are shown
without confidence intervals; the fixed formula is not parameter
matched. These are prediction diagnostics, not closed-loop results.}
\end{figure}

\subsubsection{6.2. A post-hoc support-gated
follow-up}\label{a-post-hoc-support-gated-follow-up}

A subsequent six-seed diagnostic adds learned corrections to the fixed
response estimate. The support-gated version reduces the correction when
past excitation is strong. This follow-up was selected after viewing
earlier results on the same development data.

\textbf{Table 5. Observed exploratory follow-up.} Late history means at
least four completed records.

{\def\LTcaptype{none} % do not increment counter
\begin{longtable}[]{@{}lrr@{}}
\toprule\noalign{}
Diagnostic predictor & Overall MSE & Late-history MSE \\
\midrule\noalign{}
\endhead
\bottomrule\noalign{}
\endlastfoot
Fixed formula & 7.3889e-5 & 1.6571e-5 \\
Formula plus learned residual & 5.6706e-5 & 2.3687e-5 \\
Formula plus support-gated residual & 2.4417e-5 & 1.6588e-5 \\
\end{longtable}
}

The gate mainly prevents a learned correction from damaging an already
useful mature-history estimate and improves cold-start behavior.
Late-history performance is effectively at the fixed-formula level. This
is not evidence that excitation support is calibrated uncertainty or
that the online controller should trust all high-support records.

\subsubsection{6.3. What online pilots do and do not
show}\label{what-online-pilots-do-and-do-not-show}

An existing overnight memory run reached 516 outer iterations before its
12-hour budget expired. Evaluation used four fixed trajectories: one
checkpoint achieved two successes, while later checkpoints at iterations
400, 450, and 500 achieved one success each. The run used always-on
actor control, whereas the formal baseline used a critical-phase routing
configuration. There is no matched no-memory overnight run establishing
attribution. These observations demonstrate that the training path
executes; they do not demonstrate faster convergence or a reliable 25\%
final success estimate.

Synthetic condition-shift probes further show that stale high-support
response estimates can be misleading and that delayed feedback requires
generation bookkeeping. These are software and mechanism diagnostics,
not manipulation experiments. None of the current evidence measures
saved human time, source-to-target transfer, or sustained autonomous
improvement.

\subsubsection{6.4. Confirmatory results to
populate}\label{confirmatory-results-to-populate}

\textbf{Table 6. Formal results template. Every em dash denotes an
unmeasured quantity.} Populate separately for each task and condition;
aggregate only after reporting per-task results.

{\def\LTcaptype{none} % do not increment counter
\begin{longtable}[]{@{}
  >{\raggedright\arraybackslash}p{(\linewidth - 10\tabcolsep) * \real{0.1304}}
  >{\raggedleft\arraybackslash}p{(\linewidth - 10\tabcolsep) * \real{0.1739}}
  >{\raggedleft\arraybackslash}p{(\linewidth - 10\tabcolsep) * \real{0.1739}}
  >{\raggedleft\arraybackslash}p{(\linewidth - 10\tabcolsep) * \real{0.1739}}
  >{\raggedleft\arraybackslash}p{(\linewidth - 10\tabcolsep) * \real{0.1739}}
  >{\raggedleft\arraybackslash}p{(\linewidth - 10\tabcolsep) * \real{0.1739}}@{}}
\toprule\noalign{}
\begin{minipage}[b]{\linewidth}\raggedright
Method
\end{minipage} & \begin{minipage}[b]{\linewidth}\raggedleft
Success AUC
\end{minipage} & \begin{minipage}[b]{\linewidth}\raggedleft
Final success, 95\% CI
\end{minipage} & \begin{minipage}[b]{\linewidth}\raggedleft
Valid ticks to target
\end{minipage} & \begin{minipage}[b]{\linewidth}\raggedleft
Successes/hour
\end{minipage} & \begin{minipage}[b]{\linewidth}\raggedleft
Extra source GPU-hours
\end{minipage} \\
\midrule\noalign{}
\endhead
\bottomrule\noalign{}
\endlastfoot
Original RLT & --- & --- & --- & --- & --- \\
Capacity-matched RLT & --- & --- & --- & --- & --- \\
Recurrent recent history & --- & --- & --- & --- & --- \\
Current archive attention & --- & --- & --- & --- & --- \\
Fixed response context & --- & --- & --- & --- & --- \\
Proposed memory & --- & --- & --- & --- & --- \\
eRLT-style representation & --- & --- & --- & --- & --- \\
\end{longtable}
}

A second table must report attempt-one through attempt-five success with
frozen weights and retain/clear/shuffle interventions. A third must
report related-bank, unrelated-bank, and empty-bank target performance.
These tables cannot be inferred from prediction MSE.

\subsection{7. Analysis Plan, Limitations, and Research
Decisions}\label{analysis-plan-limitations-and-research-decisions}

Three patterns would support the central explanation: the representation
predicts subsequent effects better with appropriate past interactions;
controlled memory interventions change closed-loop performance; and the
online learning advantage persists after charging pretraining and
inference cost. Improvement on only one axis warrants a narrower
conclusion. A better predictor with unchanged control value is a useful
negative result, not evidence of sample-efficient RL.

If a strong recurrent reader matches the archive, retain the simpler
method and drop claims about long-term retrieval. If fixed response
features dominate, investigate neural estimation and observability
before expanding the architecture. If source memory harms new tasks,
report the negative transfer and restrict the bank's scope. If eRLT
removes the gap, the initial bottleneck may have been current-state
representation quality rather than missing execution history. If gains
occur only with privileged contact supervision, label the method
accordingly and test robustness to noisy deployable sensors.

The principal limitations are partial observability, insufficient
excitation, contact ambiguity, finite-memory interference, and mismatch
between source and target controllers. Offline sequence prediction can
exploit time or task correlations without identifying dynamics. Frozen
visual targets may discard subtle contact effects, while changes in
viewpoint may dominate their distance. A small number of simulation
tasks does not establish universal physical knowledge. Persistent
runtime memory and exact resumed environment state require additional
systems work. Main results must specify which form of persistence was
actually implemented.

For a computer-vision submission, visual evidence must contribute to the
method and its evaluation. Merely appending a joint-response statistic
to a VLA does not establish a new visual representation result. An
honest paper can focus on interaction-conditioned visual control, but it
needs task diversity, relevant visual ablations, strong nearest-neighbor
baselines, and reproducible evidence. A simulation-only study is
possible in principle; publication readiness depends on the strength of
these results, not on the completeness of this narrative.

\subsection{8. Conclusion and Future
Work}\label{conclusion-and-future-work}

We propose to condition lightweight frozen-VLA reinforcement learning on
representations of completed physical interactions, trained for their
utility in predicting later action consequences. The design preserves a
stable visual-language coordinate system, separates bounded context
updates from parameter learning, and defines interventions that can test
whether retained experience improves control. Existing diagnostics
justify investigating the reader, but do not yet validate the proposed
method. The next decisive experiment is a matched closed-loop comparison
followed by frozen-parameter retain-versus-clear evaluation. If these
succeed, the study can expand to source-bank transfer, changed contact
conditions, joint Stage 1 training, and eventually measured operator
effort on physical robots.

\subsection{Appendix A. Implementation Status at Draft
Time}\label{appendix-a.-implementation-status-at-draft-time}

{\def\LTcaptype{none} % do not increment counter
\begin{longtable}[]{@{}
  >{\raggedright\arraybackslash}p{(\linewidth - 4\tabcolsep) * \real{0.3333}}
  >{\raggedright\arraybackslash}p{(\linewidth - 4\tabcolsep) * \real{0.3333}}
  >{\raggedright\arraybackslash}p{(\linewidth - 4\tabcolsep) * \real{0.3333}}@{}}
\toprule\noalign{}
\begin{minipage}[b]{\linewidth}\raggedright
Capability
\end{minipage} & \begin{minipage}[b]{\linewidth}\raggedright
Existing branch
\end{minipage} & \begin{minipage}[b]{\linewidth}\raggedright
Proposed work
\end{minipage} \\
\midrule\noalign{}
\endhead
\bottomrule\noalign{}
\endlastfoot
Frozen Stage 2 VLA and original RL token & Implemented & Preserve \\
Raw executed command/joint-response record & Implemented, 110 dimensions
& Preserve units and masks \\
Short store / archive & 4 recent, 32 stored, 4 retrieved & Add visual
event content and learned retrieval study \\
Current memory reader & 64-dimensional attention or fixed response & Add
consequence-supervised event encoder/reader \\
Reader update & Critic TD owns trainable reader & Freeze pretrained
reader first; online update ablation later \\
Stage 1B interaction pretraining & Not implemented in Zeva branch & Add
chronological transition dataset and auxiliary decoder \\
Same-instance retention & Optional, reset checks required & Complete
frozen-weight evaluation tooling \\
Cross-task source bank & Not implemented & Separate
provenance-controlled bank \\
Actual contact labels & Not in default event & Optional privileged
auxiliary supervision \\
Context in action expert & Not implemented & Later extension with
deployable feature contract \\
Production change detector & Not enabled & Optional robustness variant
after core result \\
\end{longtable}
}

The local inspected revision is
\texttt{f2d1cf6486d309eb61785151bc33783bc64a967f}. A newer remote
research revision, \texttt{abb4382362ae0c89f711009227cc236b6d6f8b4d},
contains delayed-feedback generation protection in diagnostics and has
not been merged into this working tree. Existing reports and raw result
files are indexed in the accompanying execution guide and evidence
manifest.

\subsection{Appendix B. Reproducibility
Record}\label{appendix-b.-reproducibility-record}

For every experiment, save the git revision and dirty patch, resolved
Hydra configuration, Stage 1 and context checkpoint hashes,
normalization statistics, simulator/controller versions,
train/validation/test identifiers, memory clear/retain policy, seeds,
source-bank provenance, actual executed ticks, intervention settings,
and hardware placement. Store evaluation per episode and per attempt. A
resume check must exceed replay cache capacity and verify retained
transition files, counts, optimizer/RNG state, and the declared
environment-memory reset behavior.

Official RLinf was checked on 3 October 2026: main
\texttt{c70606f08cdca259b8dec03d4430926b5b8fac9d}; replay PR \#1623
remained open. Upstream fixes are evaluated in a separate revision
before a campaign. A running campaign retains its fixed code version.
The document bundle contains source metadata and a snapshot of recently
updated official PRs.

\subsection*{References}\label{references}
\addcontentsline{toc}{subsection}{References}

\protect\phantomsection\label{refs}
\begin{CSLReferences}{1}{1}
\bibitem[\citeproctext]{ref-rliable}
Agarwal, Rishabh, Max Schwarzer, Pablo Samuel Castro, Aaron Courville,
and Marc G. Bellemare. 2021. \emph{Deep Reinforcement Learning at the
Edge of the Statistical Precipice}.
\url{https://arxiv.org/abs/2108.13264}.

\bibitem[\citeproctext]{ref-her}
Andrychowicz, Marcin, Filip Wolski, Alex Ray, et al. 2017.
\emph{Hindsight Experience Replay}.
\url{https://arxiv.org/abs/1707.01495}.

\bibitem[\citeproctext]{ref-pi0}
Black, Kevin, Noah Brown, Danny Driess, et al. 2024. \emph{\(\pi_0\): A
Vision-Language-Action Flow Model for General Robot Control}.
\url{https://arxiv.org/abs/2410.24164}.

\bibitem[\citeproctext]{ref-zeva}
Chen, Fu, Xin Ding, Bingjia Huang, et al. 2026. \emph{Zeva: In-Context
Causal Learning for Generalizable Embodied Manipulation}.
\url{https://arxiv.org/abs/2608.30880}.

\bibitem[\citeproctext]{ref-conrft}
Chen, Yuhui, Shuai Tian, Shugao Liu, Yingting Zhou, Haoran Li, and
Dongbin Zhao. 2025. \emph{ConRFT: A Reinforced Fine-Tuning Method for
VLA Models via Consistency Policy}.
\url{https://arxiv.org/abs/2502.05450}.

\bibitem[\citeproctext]{ref-robomme}
Dai, Yinpei, Hongze Fu, Jayjun Lee, et al. 2026. \emph{RoboMME:
Benchmarking and Understanding Memory for Robotic Generalist Policies}.
\url{https://arxiv.org/abs/2603.04639}.

\bibitem[\citeproctext]{ref-rl2}
Duan, Yan, John Schulman, Xi Chen, Peter L. Bartlett, Ilya Sutskever,
and Pieter Abbeel. 2016. \emph{RL\(^2\): Fast Reinforcement Learning via
Slow Reinforcement Learning}. \url{https://arxiv.org/abs/1611.02779}.

\bibitem[\citeproctext]{ref-td3}
Fujimoto, Scott, Herke van Hoof, and David Meger. 2018. \emph{Addressing
Function Approximation Error in Actor-Critic Methods}.
\url{https://arxiv.org/abs/1802.09477}.

\bibitem[\citeproctext]{ref-sac}
Haarnoja, Tuomas, Aurick Zhou, Pieter Abbeel, and Sergey Levine. 2018.
\emph{Soft Actor-Critic: Off-Policy Maximum Entropy Deep Reinforcement
Learning with a Stochastic Actor}.
\url{https://arxiv.org/abs/1801.01290}.

\bibitem[\citeproctext]{ref-tdmpc2}
Hansen, Nicklas, Hao Su, and Xiaolong Wang. 2023. \emph{TD-MPC2:
Scalable, Robust World Models for Continuous Control}.
\url{https://arxiv.org/abs/2310.16828}.

\bibitem[\citeproctext]{ref-imagine}
Hu, Kejia, Wentong Zhai, Bo Zhao, and Shuai Liang. 2026.
\emph{Imagine-RL: Residual-Confidence-Guided Cross-Attention for
World-Model-Augmented VLA Reinforcement Learning}.
\url{https://arxiv.org/abs/2609.24033}.

\bibitem[\citeproctext]{ref-erlt}
Huang, Dehao, Jianbang Liu, Jianpan Gao, et al. 2026. \emph{eRLT:
Efficient VLA Reinforcement Learning via Action-Relevant Token Routing}.
\url{https://arxiv.org/abs/2610.00913}.

\bibitem[\citeproctext]{ref-pi05}
Intelligence, Physical, Kevin Black, Noah Brown, et al. 2025.
\emph{\(\pi_{0.5}\): A Vision-Language-Action Model with Open-World
Generalization}. \url{https://arxiv.org/abs/2504.16054}.

\bibitem[\citeproctext]{ref-rma}
Kumar, Ashish, Zipeng Fu, Deepak Pathak, and Jitendra Malik. 2021.
\emph{RMA: Rapid Motor Adaptation for Legged Robots}.
\url{https://arxiv.org/abs/2107.04034}.

\bibitem[\citeproctext]{ref-libero}
Liu, Bo, Yifeng Zhu, Chongkai Gao, et al. 2023. \emph{LIBERO:
Benchmarking Knowledge Transfer for Lifelong Robot Learning}.
\url{https://arxiv.org/abs/2306.03310}.

\bibitem[\citeproctext]{ref-hilserl}
Luo, Jianlan, Charles Xu, Jeffrey Wu, and Sergey Levine. 2024.
\emph{Precise and Dexterous Robotic Manipulation via Human-in-the-Loop
Reinforcement Learning}. \url{https://arxiv.org/abs/2410.21845}.

\bibitem[\citeproctext]{ref-pearl}
Rakelly, Kate, Aurick Zhou, Deirdre Quillen, Chelsea Finn, and Sergey
Levine. 2019. \emph{Efficient Off-Policy Meta-Reinforcement Learning via
Probabilistic Context Variables}.
\url{https://arxiv.org/abs/1903.08254}.

\bibitem[\citeproctext]{ref-spr}
Schwarzer, Max, Ankesh Anand, Rishab Goel, R Devon Hjelm, Aaron
Courville, and Philip Bachman. 2020. \emph{Data-Efficient Reinforcement
Learning with Self-Predictive Representations}.
\url{https://arxiv.org/abs/2007.05929}.

\bibitem[\citeproctext]{ref-memoryvla}
Shi, Hao, Bin Xie, Yingfei Liu, et al. 2025. \emph{MemoryVLA:
Perceptual-Cognitive Memory in Vision-Language-Action Models for Robotic
Manipulation}. \url{https://arxiv.org/abs/2508.19236}.

\bibitem[\citeproctext]{ref-maniskill}
Tao, Stone, Fanbo Xiang, Arth Shukla, et al. 2024. \emph{ManiSkill3: GPU
Parallelized Robotics Simulation and Rendering for Generalizable
Embodied AI}. \url{https://arxiv.org/abs/2410.00425}.

\bibitem[\citeproctext]{ref-dsrl}
Wagenmaker, Andrew, Mitsuhiko Nakamoto, Yunchu Zhang, et al. 2025.
\emph{Steering Your Diffusion Policy with Latent Space Reinforcement
Learning}. \url{https://arxiv.org/abs/2506.15799}.

\bibitem[\citeproctext]{ref-rlt}
Xu, Charles, Jost Tobias Springenberg, Michael Equi, et al. 2026.
\emph{RL Token: Bootstrapping Online RL with Vision-Language-Action
Models}. \url{https://arxiv.org/abs/2604.23073}.

\bibitem[\citeproctext]{ref-uposi}
Yu, Wenhao, Jie Tan, C. Karen Liu, and Greg Turk. 2017. \emph{Preparing
for the Unknown: Learning a Universal Policy with Online System
Identification}. \url{https://arxiv.org/abs/1702.02453}.

\bibitem[\citeproctext]{ref-flare}
Zheng, Ruijie, Jing Wang, Scott Reed, et al. 2025. \emph{FLARE: Robot
Learning with Implicit World Modeling}.
\url{https://arxiv.org/abs/2505.15659}.

\bibitem[\citeproctext]{ref-varibad}
Zintgraf, Luisa, Kyriacos Shiarlis, Maximilian Igl, et al. 2019.
\emph{VariBAD: A Very Good Method for Bayes-Adaptive Deep RL via
Meta-Learning}. \url{https://arxiv.org/abs/1910.08348}.

\end{CSLReferences}

\end{document}
